\documentclass[a4paper]{arrowhead}

\usepackage[yyyymmdd]{datetime}
\usepackage{etoolbox}
\usepackage[T1]{fontenc}
\usepackage[utf8]{inputenc}
\usepackage{multirow}
\usepackage{longtable}
\usepackage{amssymb}
\usepackage{float}
\usepackage{tikz}
\usetikzlibrary{positioning,arrows.meta,fit,calc,shapes.misc,shapes.geometric}

\renewcommand{\dateseparator}{-}
\setlength{\emergencystretch}{3em}
\lstset{basicstyle=\ttfamily\footnotesize,breaklines=true,columns=fullflexible,frame=single,rulecolor=\color{gray!50},backgroundcolor=\color{gray!6}}
\tikzset{
  ahsys/.style={draw=ArrowheadBlue,very thick,fill=ArrowheadBlue!6,rounded corners=2pt,minimum width=3.2cm,minimum height=1.1cm,align=center,font=\small\bfseries},
  ahext/.style={draw=ArrowheadGrey,thick,dashed,fill=gray!5,rounded corners=2pt,minimum width=3.2cm,minimum height=1.0cm,align=center,font=\small},
  ahsrv/.style={draw=ArrowheadPurple,thick,fill=ArrowheadPurple!8,rounded corners=6pt,minimum height=0.55cm,inner xsep=4pt,align=center,font=\scriptsize\ttfamily},
  ahcsm/.style={draw=ArrowheadTeal,thick,fill=ArrowheadTeal!8,rounded corners=6pt,minimum height=0.55cm,inner xsep=4pt,align=center,font=\scriptsize\ttfamily},
  ahdb/.style={draw=ArrowheadOrange,thick,fill=ArrowheadOrange!8,cylinder,shape border rotate=90,aspect=0.25,minimum height=0.9cm,align=center,font=\scriptsize},
  ahflow/.style={-{Stealth[length=2.2mm]},thick,ArrowheadGrey},
  ahlabel/.style={font=\scriptsize,fill=white,inner sep=1pt,align=center},
}

%% Special references
\newcommand{\fref}[1]{{\textcolor{ArrowheadBlue}{\hyperref[sec:functions:#1]{#1}}}}
\newcommand{\mref}[1]{{\textcolor{ArrowheadPurple}{\hyperref[sec:model:#1]{#1}}}}
\newcommand{\sref}[1]{{\textcolor{ArrowheadPurple}{#1}}}
\newrobustcmd\msubsection[2]{
  \addtocounter{subsection}{1}
  \addcontentsline{toc}{subsection}{\protect\numberline{\thesubsection}#1 \textcolor{ArrowheadPurple}{#2}}
  \subsection*{\thesubsection\quad#1 \textcolor{ArrowheadPurple}{#2}}
  \label{sec:model:#2}
}

\begin{document}

%% Arrowhead Document Properties
\ArrowheadTitle{WorkflowManager}

\ArrowheadType{System Design Description}
\ArrowheadTypeShort{SysDD}
\ArrowheadVersion{5.4.0}
\ArrowheadDate{2026-11-24}
\ArrowheadAuthor{Documentation maintenance}
\ArrowheadStatus{REVIEW DRAFT}
\ArrowheadContact{Repository maintainers}
\ArrowheadFooter{\href{www.arrowhead.eu}{www.arrowhead.eu}}
\ArrowheadSetup
%%

%% Front Page
\begin{center}
  \vspace*{1cm}
  \huge{\arrowtitle}

  \vspace*{0.2cm}
  \LARGE{\arrowtype}
  \vspace*{1cm}

  
  \vspace*{\fill}

  

  \vspace*{1cm}
  \vspace*{\fill}

  % Front Page Abstract
  \begin{abstract}
    This document describes how the WorkflowManager is implemented in the Java reference implementation: its internal components, the realisation of its functional and non-functional properties, its configuration and the implementation status of the other languages. Documentation edition 5.4.0.
  \end{abstract}

  \vspace*{1cm}
\end{center}

\newpage
%%

%% Table of Contents
\tableofcontents
\newpage
%%

\section{Overview}
\label{sec:overview}
This document describes the design of the WorkflowManager, which keeps the books of workflows: it holds recipes and sessions, validates recipes, advances sessions with the transition rules of the formal model, dispatches steps to WorkflowExecutors and reports to the product. The abstract system is described in the
SysD of the same name; its produced services are specified by SDs and IDDs in this documentation set.

\begin{figure}[ht!]
  \centering
  \resizebox{\textwidth}{!}{
\begin{tikzpicture}[node distance=4mm]
\node[ahsys,minimum height=5.25cm,minimum width=3.6cm] (sys) at (0,0) {WorkflowManager\\[-1pt]{\normalfont\scriptsize $\ll$system$\gg$}};
\node[ahsrv,anchor=west] (p0) at (3.4,2.25) {workflowRecipe};
\draw[thick,ArrowheadPurple] (1.8,2.25) -- (p0.west);
\fill[ArrowheadPurple] (1.8,2.25) circle (1.4pt);
\node[ahsrv,anchor=west] (p1) at (3.4,1.50) {workflowInstance};
\draw[thick,ArrowheadPurple] (1.8,1.50) -- (p1.west);
\fill[ArrowheadPurple] (1.8,1.50) circle (1.4pt);
\node[ahsrv,anchor=west] (p2) at (3.4,0.75) {workflowWorkItem};
\draw[thick,ArrowheadPurple] (1.8,0.75) -- (p2.west);
\fill[ArrowheadPurple] (1.8,0.75) circle (1.4pt);
\node[ahsrv,anchor=west] (p3) at (3.4,0.00) {workflowContext};
\draw[thick,ArrowheadPurple] (1.8,0.00) -- (p3.west);
\fill[ArrowheadPurple] (1.8,0.00) circle (1.4pt);
\node[ahsrv,anchor=west] (p4) at (3.4,-0.75) {workflowLease};
\draw[thick,ArrowheadPurple] (1.8,-0.75) -- (p4.west);
\fill[ArrowheadPurple] (1.8,-0.75) circle (1.4pt);
\node[ahsrv,anchor=west] (p5) at (3.4,-1.50) {workflowManagement};
\draw[thick,ArrowheadPurple] (1.8,-1.50) -- (p5.west);
\fill[ArrowheadPurple] (1.8,-1.50) circle (1.4pt);
\node[ahsrv,anchor=west] (p6) at (3.4,-2.25) {generalManagement};
\draw[thick,ArrowheadPurple] (1.8,-2.25) -- (p6.west);
\fill[ArrowheadPurple] (1.8,-2.25) circle (1.4pt);
\node[ahcsm,anchor=east] (c0) at (-3.4,2.25) {serviceDiscovery\\[-2pt]{\tiny\rmfamily ServiceRegistry}};
\draw[thick,ArrowheadTeal,-{Stealth[length=1.6mm]}] (c0.east) -- (-1.8,2.25);
\node[ahcsm,anchor=east] (c1) at (-3.4,1.50) {identity\\[-2pt]{\tiny\rmfamily Authentication}};
\draw[thick,ArrowheadTeal,-{Stealth[length=1.6mm]}] (c1.east) -- (-1.8,1.50);
\node[ahcsm,anchor=east] (c2) at (-3.4,0.75) {authorization\\[-2pt]{\tiny\rmfamily ConsumerAuthorization}};
\draw[thick,ArrowheadTeal,-{Stealth[length=1.6mm]}] (c2.east) -- (-1.8,0.75);
\node[ahcsm,anchor=east] (c3) at (-3.4,0.00) {workflowExecution\\[-2pt]{\tiny\rmfamily WorkflowExecutor}};
\draw[thick,ArrowheadTeal,-{Stealth[length=1.6mm]}] (c3.east) -- (-1.8,0.00);
\node[ahcsm,anchor=east] (c4) at (-3.4,-0.75) {workflowProduct\\[-2pt]{\tiny\rmfamily Intelligent Product}};
\draw[thick,ArrowheadTeal,-{Stealth[length=1.6mm]}] (c4.east) -- (-1.8,-0.75);
\node[ahcsm,anchor=east] (c5) at (-3.4,-1.50) {sysMonitor (AMQP)\\[-2pt]{\tiny\rmfamily SysMonitor broker}};
\draw[thick,ArrowheadTeal,-{Stealth[length=1.6mm]}] (c5.east) -- (-1.8,-1.50);
\node[font=\scriptsize\bfseries,ArrowheadPurple] at (5.0,3.15) {produced};
\node[font=\scriptsize\bfseries,ArrowheadTeal] at (-5.0,3.15) {consumed};
\end{tikzpicture}
  }
  \caption{Block view of the WorkflowManager with its produced and consumed services.}
  \label{fig:block}
\end{figure}


\newpage
\section{Implementation}
\label{sec:implementation}

\subsection{Implementation language and tools}
Java~21 and Spring Boot~4 (Spring Framework~7, Spring Security~7), Spring Data JPA with Hibernate, Jackson for JSON,
Springdoc for the OpenAPI description, built with Maven~3.9 (in the \texttt{maven:3.\allowbreak{}9-\allowbreak{}eclipse-\allowbreak{}temurin-\allowbreak{}21} image). The
system builds on the vendored Arrowhead~5.4.0 common library (\texttt{vendor/\allowbreak{}common-\allowbreak{}utils},
\texttt{vendor/\allowbreak{}data-\allowbreak{}transfer-\allowbreak{}objects}): system registration, identity, TLS, management filters and HTTP consumption
of other services. The source is module \texttt{workflow-\allowbreak{}manager} of \texttt{Eclipse-\allowbreak{}Arrowhead-\allowbreak{}java/\allowbreak{}ah5-\allowbreak{}Workflow-\allowbreak{}java}; the
container image runs the Spring Boot jar on Eclipse Temurin~21.

\subsection{Internal structure}

{\small
\begin{longtable}{|p{6.2cm}|p{9.6cm}|}

\hline
\rowcolor{gray!33} Component & Responsibility \\ \hline
\endhead
{\footnotesize \texttt{Workflow\allowbreak{}Manager\allowbreak{}System\allowbreak{}Info}} & Registration: system model, the seven services and their HTTP operations. \\ \hline
{\footnotesize \texttt{api.\allowbreak{}http.\allowbreak{}RecipeAPI,\allowbreak{} InstanceAPI,\allowbreak{} WorkItemAPI,\allowbreak{} ContextAPI,\allowbreak{} LeaseAPI,\allowbreak{} WorkflowManagementAPI,\allowbreak{} GeneralManagementAPI}} & HTTP controllers of the services; DTOs in \texttt{api.dto}; errors mapped by \texttt{Workflow\allowbreak{}Exception\allowbreak{}Advice} (409) and the common handler. \\ \hline
{\footnotesize \texttt{engine.\allowbreak{}WorkflowEngine}} & Drives sessions: load the snapshot, apply rules (\texttt{Workflow\allowbreak{}Session} of \texttt{workflow-\allowbreak{}model}), dispatch, handle claims and reports, watchdog for claim timeouts, step timeouts and retry backoffs, resume after restart. One lock per session. \\ \hline
{\footnotesize \texttt{service.\allowbreak{}RecipeService,\allowbreak{} CatalogueCheck}} & Catalogue: parse (unknown fields rejected), validate (\texttt{Recipe\allowbreak{}Validator} plus the \texttt{RecipeCheck} SPI with C01/C02), store canonical JSON, retire, compile. \\ \hline
{\footnotesize \texttt{service.\allowbreak{}ExecutorDirectory,\allowbreak{} HttpExecutorGateway}} & Executor set $WE$ from the ServiceRegistry (\texttt{workflow\allowbreak{}Execution} metadata) and static configuration; push calls to executors. \\ \hline
{\footnotesize \texttt{service.\allowbreak{}LeaseService}} & FIFO leases with tokens, renewal, expiry and session release. \\ \hline
{\footnotesize \texttt{service.\allowbreak{}ProductGateway}} & Consumer of \texttt{workflow\allowbreak{}Product}. \\ \hline
{\footnotesize \texttt{service.\allowbreak{}WorklogService,\allowbreak{} jpa.\allowbreak{}*}} & Worklog and persistence: Spring Data JPA on the embedded H2 database; the system applies \texttt{database/\allowbreak{}schema.\allowbreak{}sql} at start and Hibernate only validates it. \\ \hline
{\footnotesize \texttt{security.\allowbreak{}ServiceAuthorizationFilter,\allowbreak{} ServiceAuthorizationService}} & Provider-side ConsumerAuthorization check per operation. \\ \hline
{\footnotesize \texttt{workflow-\allowbreak{}model (library)}} & Recipe, codec, graph, validator, session rules, guard expressions, input composer, plan compiler, wire records: pure Java and Jackson, no Spring. \\ \hline
\end{longtable}
}


\subsection{Functional properties implementation}
A request that changes a session (start, report, abort, claim timeout) takes the session's lock, loads the
snapshot of $z$, applies one or more rules of \texttt{Workflow\allowbreak{}Session}, saves the snapshot (optimistic revision) and
appends the transitions to the trace and worklog. Outbound calls (dispatch to executors, aborts, product calls) are
collected while the lock is held and executed after it is released, on the worker pool. Executor selection in push mode
filters $WE$ by $\eta(s)$ and $\theta(s)$ (service names compared case-insensitively without separators), excludes
executors whose work item of the step failed, and picks the least loaded. After a \texttt{BIND\_FAILED} report or an
unreachable executor the next executor is tried; this is S-Bind with another candidate and records no extension. When
none is left S-FailBind fires. After an \texttt{ERROR} report, while the retry policy allows another attempt, x-rebind
excludes the executor and dispatches the step again, whereas x-retry keeps the attempt in state \texttt{RETRY\_WAIT} for
the backoff and then dispatches to the same executor, if it is still a candidate. The backoff is ended by a delayed task
or, after a restart, by the watchdog, which runs every second. In pull mode a retried step is published again and any
suitable executor may claim it. Guards are evaluated by
\texttt{Guard\allowbreak{}Expression} from $\rho$ and $\xi$ when a step completes; a step whose incoming transitions were all not
taken is skipped (x-skip).

\subsection{Non functional properties implementation}
\subsubsection{Security}
Spring Security filter chains of the common library: TLS client authentication, identity-token verification against
the Authentication system (outsourced authentication), the management filter and the workflow systems'
\texttt{Service\allowbreak{}Authorization\allowbreak{}Filter}, which maps method and path to service definition and operation and asks
ConsumerAuthorization (\texttt{authorization.\allowbreak{}verify}). Object-level checks (holder, owner, configured manager) are made in
the service logic and give 403.
\subsubsection{Power management}
No specific power management. Thread pools are bounded; the WorkflowManager's watchdog and the executor's pull loop are
the only periodic activities.
\subsubsection{Internal monitoring}
Spring Boot health endpoint (\texttt{Workflow\allowbreak{}Manager\allowbreak{}Health\allowbreak{}Indicator}),
the common-library log table exposed through \texttt{general\allowbreak{}Management.\allowbreak{}get-\allowbreak{}log}, and optional sysMonitor telemetry over
AMQP (\texttt{sysmonitor.\allowbreak{}enabled}). The worklog records every session event.
\subsubsection{Configuration}

{\small
\begin{longtable}{|p{5.6cm}|p{3.4cm}|p{6.8cm}|}

\hline
\rowcolor{gray!33} Property & Default & Meaning \\ \hline
\endhead
\texttt{server.\allowbreak{}port} & 8482 & Arrowhead connector port. \\ \hline
\texttt{workflow.\allowbreak{}workstation.\allowbreak{}id} & (empty) & Workstation identifier sent to products. \\ \hline
\texttt{workflow.\allowbreak{}dispatch.\allowbreak{}default} & push & Dispatch mode when neither the session nor the recipe gives one. \\ \hline
\texttt{workflow.\allowbreak{}executors} & (empty) & Statically configured executors (development, or executors outside the ServiceRegistry). \\ \hline
\texttt{workflow.\allowbreak{}executor.\allowbreak{}discovery.\allowbreak{}cache-\allowbreak{}seconds} & 10 & How long discovered executors are cached. \\ \hline
\texttt{workflow.\allowbreak{}pull.\allowbreak{}claim-\allowbreak{}timeout} & 300 & Seconds a published (pull) work item may wait for a claim before S-FailBind. \\ \hline
\texttt{workflow.\allowbreak{}worker.\allowbreak{}threads} & 8 & Dispatch worker threads. \\ \hline
\texttt{workflow.\allowbreak{}lease.\allowbreak{}default-\allowbreak{}duration} & 120 & Default lease duration in seconds. \\ \hline
\texttt{workflow.\allowbreak{}lease.\allowbreak{}release-\allowbreak{}on-\allowbreak{}failure} & false & Compensation: release a failed session's leases. \\ \hline
\texttt{workflow.\allowbreak{}validator.\allowbreak{}catalogue-\allowbreak{}check} & true & Run C02 (executor availability) against the ServiceRegistry. \\ \hline
\texttt{workflow.\allowbreak{}validator.\allowbreak{}catalogue-\allowbreak{}check.\allowbreak{}services} & true & Run C01 (service registered). The secure deployment sets it to false, because under RESTRICTED discovery the lookup hides providers without \texttt{unrestricted\allowbreak{}Discovery}. \\ \hline
\texttt{workflow.\allowbreak{}standalone} & false & Development mode without core systems. \\ \hline
\texttt{service.\allowbreak{}authorization.\allowbreak{}enabled} & false (deployment: true) & Provider-side ConsumerAuthorization check. \\ \hline
\texttt{management.\allowbreak{}policy} & (deployment: sysop-only) & Protection of the management services. \\ \hline
\texttt{DATABASE\_\allowbreak{}DIRECTORY} & ./data (deployment: /data) & Directory of the embedded database; a volume of this system only. \\ \hline
\texttt{spring.\allowbreak{}datasource.\allowbreak{}password} & (secret ah.db-password) & Opens the database file. Set before the first start and kept: a changed value locks the system out of its data. \\ \hline
\texttt{sysmonitor.\allowbreak{}enabled} & true (deployment: false) & Publish sysMonitor telemetry over AMQP; the deployment enables it with \texttt{SYSMONITOR\_\allowbreak{}ENABLED}. \\ \hline
\end{longtable}
}


\newpage
\section{Services}
\label{sec:services}
\subsection{Produced services}
All produced services use interface template \texttt{generic\_\allowbreak{}https}, policy \texttt{CERT\_AUTH} and version 1.0.0 and are
registered with every operation's method and path.

{\small
\begin{longtable}{|p{3.8cm}|p{6.2cm}|p{5.8cm}|}

\hline
\rowcolor{gray!33} Service & Operations & Documents \\ \hline
\endhead
\texttt{workflow\allowbreak{}Recipe} & validate, register, query, get, retire, compile & \texttt{Workflow\allowbreak{}Manager-\allowbreak{}v5.\allowbreak{}4.\allowbreak{}0-\allowbreak{}workflow\allowbreak{}Recipe-\allowbreak{}SD}, \texttt{-HTTP-IDD} \\ \hline
\texttt{workflow\allowbreak{}Instance} & start, abort, query, get-status, get-trace & \texttt{Workflow\allowbreak{}Manager-\allowbreak{}v5.\allowbreak{}4.\allowbreak{}0-\allowbreak{}workflow\allowbreak{}Instance-\allowbreak{}SD}, \texttt{-HTTP-IDD} \\ \hline
\texttt{workflow\allowbreak{}Work\allowbreak{}Item} & claim, report & \texttt{Workflow\allowbreak{}Manager-\allowbreak{}v5.\allowbreak{}4.\allowbreak{}0-\allowbreak{}workflow\allowbreak{}Work\allowbreak{}Item-\allowbreak{}SD}, \texttt{-HTTP-IDD} \\ \hline
\texttt{workflow\allowbreak{}Context} & get-parameters, get-results, put-result & \texttt{Workflow\allowbreak{}Manager-\allowbreak{}v5.\allowbreak{}4.\allowbreak{}0-\allowbreak{}workflow\allowbreak{}Context-\allowbreak{}SD}, \texttt{-HTTP-IDD} \\ \hline
\texttt{workflowLease} & reserve, poll, renew, release & \texttt{Workflow\allowbreak{}Manager-\allowbreak{}v5.\allowbreak{}4.\allowbreak{}0-\allowbreak{}workflow\allowbreak{}Lease-\allowbreak{}SD}, \texttt{-HTTP-IDD} \\ \hline
\texttt{workflow\allowbreak{}Management} & query-worklog, query-executors, query-work-items & \texttt{Workflow\allowbreak{}Manager-\allowbreak{}v5.\allowbreak{}4.\allowbreak{}0-\allowbreak{}workflow\allowbreak{}Management-\allowbreak{}SD}, \texttt{-HTTP-IDD} \\ \hline
\texttt{general\allowbreak{}Management} & get-log, get-config & \texttt{Workflow-\allowbreak{}v5.\allowbreak{}4.\allowbreak{}0-\allowbreak{}general\allowbreak{}Management-\allowbreak{}SD}, \texttt{-HTTP-IDD} \\ \hline
\end{longtable}
}

\subsection{Consumed services}

{\small
\begin{longtable}{|p{3.4cm}|p{3.2cm}|p{3.6cm}|p{5.6cm}|}

\hline
\rowcolor{gray!33} Service & Provider & Operations & Purpose \\ \hline
\endhead
\texttt{service\allowbreak{}Discovery} & ServiceRegistry & register, revoke, lookup & Registration; discovery of executors (\texttt{workflow\allowbreak{}Execution} with its metadata) and the catalogue check C01. \\ \hline
\texttt{identity} & Authentication & login, logout, identity-verify & Own identity token; verification of callers. \\ \hline
\texttt{authorization} & ConsumerAuthorization & verify & Provider-side authorization of every service operation. \\ \hline
\texttt{workflow\allowbreak{}Execution} & WorkflowExecutor & start, abort & Push dispatch and abort of work items. \\ \hline
\texttt{workflow\allowbreak{}Product} & Intelligent Product & retrieve-process-step, store-process-step-result & Product-driven sessions. Contract: \texttt{Workflow-\allowbreak{}v5.\allowbreak{}4.\allowbreak{}0-\allowbreak{}workflow\allowbreak{}Product-\allowbreak{}SD}. \\ \hline
\texttt{sysMonitor (AMQP)} & SysMonitor broker & publish & Health and metrics when \texttt{sysmonitor.\allowbreak{}enabled}. \\ \hline
\end{longtable}
}


\section{Verification}
\label{sec:verification}
Unit tests: 35 \texttt{workflow-\allowbreak{}model} tests (rules, validator, guards, codec, compiler) and 15 WorkflowManager tests
(Spring, with the system's own schema on the embedded database and Hibernate validating it), passed on 2026-11-24. The
smoke test (\texttt{verification/\allowbreak{}smoke}: one WorkflowManager with its embedded database, three executors, simulated plant)
passed 18/18 checks on the same day, including three concurrent forklift sessions without overlapping tool
use and lease expiry after a blocked AGV. The cloud test (\texttt{verification/\allowbreak{}cloud/\allowbreak{}run-\allowbreak{}cloud-\allowbreak{}test.\allowbreak{}sh}) against the
Java test cloud with RESTRICTED discovery, identity tokens and ConsumerAuthorization passed 49/49 checks on 2026-11-21,
before the database was embedded; it has not been repeated since: registration,
executor discovery, 401/403 cases, DSO binding with the executors' certificates, the transition counts 11, 33 and 9,
pull mode, owner-only abort, compile, guarded recipes and sysop-only management.


\section{Implementation variants}
\label{sec:variants}
\begin{itemize}
\item \textbf{Java (reference implementation).} \texttt{Eclipse-\allowbreak{}Arrowhead-\allowbreak{}java/\allowbreak{}ah5-\allowbreak{}Workflow-\allowbreak{}java/\allowbreak{}workflow-\allowbreak{}manager}, Java 21, Spring
Boot, built with Maven together with the shared \texttt{workflow-\allowbreak{}model} module and the vendored Arrowhead~5.4.0 common
library.
\item \textbf{Go.} Not implemented. No workflow system exists in \texttt{Eclipse-\allowbreak{}Arrowhead-\allowbreak{}go}.
\item \textbf{Rust.} Not implemented. No workflow system exists in \texttt{Eclipse-\allowbreak{}Arrowhead-\allowbreak{}rust}.
\end{itemize}
The workflow systems are therefore not yet interoperable across the three implementations. A Go or Rust port must
reproduce the IDDs of this set exactly and, for the WorkflowManager, the behaviour of \texttt{workflow-\allowbreak{}model}: the
recipe schema and decoder, the validator's findings and verdicts, the guard grammar and the transition rules. The
architecture prescribes conformance vectors (recipe in, verdict and transition trace out) to be run by all three
implementations; they are not written yet.


\newpage
\bibliographystyle{IEEEtran}
\bibliography{bibliography}

\newpage

\section{Revision History}
\subsection{Amendments}

\noindent\begin{tabularx}{\textwidth}{| p{1cm} | p{2.2cm} | p{1.6cm} | X | p{3.6cm} |} \hline
\rowcolor{gray!33} No. & Date & Version & Subject of Amendments & Author \\ \hline
1 & 2026-11-23 & \arrowversion & Initial edition: system design description of the WorkflowManager, written from the Java reference implementation (ah5-Workflow-java). & Documentation maintenance \\ \hline
2 & 2026-11-24 & \arrowversion & Embedded database instead of a separate database server; extension x-retry, extension names and retry states as implemented; claim timeout, defaults and verification status. & Documentation maintenance \\ \hline
\end{tabularx}

\subsection{Quality Assurance}

\noindent\begin{tabularx}{\textwidth}{| p{1cm} | p{2.2cm} | p{1.6cm} | X |} \hline
\rowcolor{gray!33} No. & Date & Version & Approved by \\ \hline
1 & & \arrowversion & Pending review \\ \hline
\end{tabularx}

\end{document}
